% main.tex — Relazione breve: sistema di telemedicina supervised_telemedicina
% Standalone: compila con make -> ../sintesi_sistema.pdf
\documentclass[11pt,a4paper]{article}

\usepackage[utf8]{inputenc}
\usepackage[T1]{fontenc}
\usepackage[italian]{babel}
\usepackage[a4paper,margin=2.5cm]{geometry}
\usepackage{amsmath,amssymb}
\usepackage{graphicx}
\usepackage{booktabs}
\usepackage{hyperref}
\usepackage{tikz}
\usetikzlibrary{positioning,arrows.meta,shapes.geometric,calc}
\usepackage{listings}
\usepackage{xcolor}

\hypersetup{colorlinks=true,linkcolor=blue!60!black,urlcolor=blue!60!black,citecolor=green!40!black,
  pdftitle={Sistema di telemedicina: sintesi del sistema}}

\definecolor{codebg}{RGB}{245,245,245}
\lstdefinestyle{compact}{language=Python,backgroundcolor=\color{codebg},
  basicstyle=\ttfamily\footnotesize,breaklines=true,frame=single,
  rulecolor=\color{black!20},numbers=none,tabsize=4,aboveskip=4pt,belowskip=4pt,
  literate={à}{{\`a}}1 {è}{{\`e}}1 {é}{{\'e}}1 {ì}{{\`i}}1 {ò}{{\`o}}1 {ù}{{\`u}}1}
\lstset{style=compact}
\newcommand{\codice}[1]{\texttt{#1}}
\newcommand{\file}[1]{\texttt{#1}}

\begin{document}

% =====================================================================
% Copertina esecutiva
\begin{center}
{\Large\bf Telemedicina supervisionata: MLP da 323 parametri con safety gate\par}
{\large accuracy 0.9851 sul test, recall sui critici 1.0 per costruzione\par}
\vspace{4pt}
{\normalsize Sintesi del progetto \file{supervised\_telemedicina} \quad\textbullet\quad seed 41 \quad\textbullet\quad \today\par}
\end{center}

\begin{quote}\small
\textbf{Abstract.} Il task è classificazione supervisionata di sei parametri
vitali in tre classi di rischio (\codice{basso}/\codice{medio}/\codice{alto}),
più risposta \codice{errore} per input invalidi. Un teacher rule-based
(\file{safety/safety\_rules.py}) etichetta offline 40\,000 + 10\,000 + 20\,000
campioni sintetici (seed 41/42/43); un MLP numpy da 323 parametri (6~$\to$~32~$\to$~3;
griglia: hidden 32, lr 0.05, batch 32) ne distilla le regole. A runtime il
safety gate ha precedenza assoluta e la soglia di incertezza 0.6 attiva il
fallback \codice{max(regole, MLP)}; ogni esito è persistito su SQLite, con
notifica al medico per critici e invalidi. Sul test congelato (20\,000):
accuracy 0.9851, kappa 0.9775, recall \codice{alto} 0.9958 (i 30 casi persi
dall'MLP sono coperti dal gate: sistema 1.0); senza buffer zone 0.9324.
Limite: 0.9851 misura fedeltà a regole sintetiche, non validità clinica.
\end{quote}

\begin{table}[h]
\centering
\caption{Risultati principali (da \file{data/models/report.json}, seed 41, test 20\,000).}
\label{tab:ris}
\begin{tabular}{lcl}
\toprule
\textbf{Misura} & \textbf{Valore} & \textbf{Note} \\
\midrule
Accuracy test & 0.9851 & 19\,702/20\,000; kappa 0.9775 \\
Recall classe \codice{alto} & 0.9958 $\to$ 1.0 & 30 persi dall'MLP, coperti dal gate \\
Accuracy senza buffer & 0.9324 & casi limite inclusi; kappa 0.8959 \\
Baseline rule-based & 1.0 & accordo del teacher sulle proprie label \\
\bottomrule
\end{tabular}
\end{table}

\begin{figure}[h]
\centering
\resizebox{\linewidth}{!}{%
\begin{tikzpicture}[
  box/.style={rectangle,draw,rounded corners,minimum width=2.2cm,minimum height=0.9cm,align=center,font=\footnotesize},
  dia/.style={diamond,draw,aspect=2,align=center,font=\footnotesize,inner sep=2pt},
  arr/.style={-{Stealth},thick,shorten >=2pt,shorten <=2pt},
  lab/.style={font=\scriptsize,align=center}]
  \node[box] (v) {6 parametri\\vitali};
  \node[box,right=0.6cm of v] (g) {safety gate\\rule-based};
  \node[dia,right=0.6cm of g] (d) {critico\\o invalido?};
  \node[box,right=0.6cm of d] (m) {MLP numpy\\conf.\ $\ge$ 0.6};
  \node[box,right=0.6cm of m] (o) {messaggio\\paziente};
  \node[box,right=0.6cm of o] (db) {SQLite +\\alert medico};
  \draw[arr] (v) -- (g);
  \draw[arr] (g) -- (d);
  \draw[arr] (d) -- node[lab,above]{no} (m);
  \draw[arr] (m) -- node[lab,above]{s\`i} (o);
  \draw[arr] (o) -- (db);
  \draw[arr,rounded corners] (d.north) -- ++(0,0.55) -| node[lab,pos=0.25,above]{s\`i: bypass MLP} (o.north);
  \node[lab,below=0.15cm of m] {sotto 0.6: fallback max(regole, MLP)};
\end{tikzpicture}%
}
\caption{Schema del sistema: il gate decide
prima dell'MLP; critici e invalidi lo bypassano.}
\label{fig:schema}
\end{figure}

% =====================================================================
\section{Architettura e pipeline}\label{sec:architettura}

A runtime ogni analisi segue un unico percorso:
\codice{CLI} $\to$ \codice{AnalysisService.analizza()} $\to$
\codice{SupervisedAgent} (safety gate $\to$ MLP) $\to$ SQLite + notifiche.
La CLI (\file{main.py}) si limita a chiamare il service: un eventuale canale
futuro (API, webapp) riusa persistenza e notifiche passando dallo stesso
punto. Offline e runtime si legano in cinque passi: (1)~il teacher rule-based
etichetta i dati; (2)~il dataset sintetico congelato (40/10/20k, seed
41/42/43) non viene mai rigenerato; (3)~la griglia prova 6 configurazioni MLP
e trattiene quella con loss di validazione minima, con riaddestramento finale
e scaler salvato nell'artifact; (4)~l'agente applica gate $\to$ MLP $\to$ soglia di incertezza $\to$
\codice{max(regole, MLP)}; (5)~\codice{AnalysisService.analizza()} valida l'input, interroga
l'agente, scrive su SQLite e, per critici e invalidi, registra la notifica in
\codice{notifiche} e la stampa su stderr. Dettagli in Sez.~\ref{sec:dati}--\ref{sec:sicurezza}.

% =====================================================================
\section{Risultati}\label{sec:risultati}

Tutti i numeri vengono da \file{data/models/report.json} (seed 41). Sul test
congelato (20\,000 campioni) l'MLP raggiunge accuracy \textbf{0.9851}
(19\,702 corrette) e kappa di Cohen \textbf{0.9775}; per la classe
\codice{alto} precision 0.9863 e recall \textbf{0.9958}: su 7\,100 veri
\codice{alto} ne riconosce 7\,070 e ne perde 30, coperti dal gate
(Sez.~\ref{sec:sicurezza}). Gli errori cadono quasi sempre tra classi vicine:
nessun \codice{basso} è predetto \codice{alto} e viceversa
(Tab.~\ref{tab:conf}).

\begin{table}[htbp]
\centering
\caption{Matrice di confusione sul test (20\,000; righe = vero, colonne = predetto).}
\label{tab:conf}
\begin{tabular}{lrrr}
\toprule
 & \textbf{pred.\ basso} & \textbf{pred.\ medio} & \textbf{pred.\ alto} \\
\midrule
vero \codice{basso} & 7\,113 & 107 & 0 \\
vero \codice{medio} & 63 & 5\,519 & 98 \\
vero \codice{alto} & 0 & 30 & 7\,070 \\
\bottomrule
\end{tabular}
\end{table}

Con $\kappa=0.9775$ l'accordo col teacher è quasi perfetto. Il test senza
buffer (campioni vicini alle soglie inclusi) scende a \textbf{0.9324} (kappa
0.8959): il numero ufficiale è valido ma ottimistico sui casi limite. I
\textbf{298 errori} si addensano ai confini (distanza media dalle soglie
0.524, massima 0.681): dove il teacher è netto l'MLP impara bene, dove la
classe cambia a gradino i confini curvi del modello sbagliano di una classe.
La baseline rule-based vale 1.0 sulle proprie label: 0.9851 misura fedeltà
alle regole sintetiche, non accuratezza su pazienti reali. Il codice è
coperto da \textbf{139 test} (unittest, comandi in Sez.~\ref{sec:uso}), più 14
di regressione legacy mai importati dal runtime.

% =====================================================================
\section{Dati e teacher rule-based}\label{sec:dati}

Le sei feature (\codice{FEATURE\_ORDER} in \file{safety\_rules.py}:
pressioni sistolica/diastolica, frequenza cardiaca, temperatura, saturazione,
glicemia; unità e range in Tab.~\ref{tab:feature}) seguono standardizzazione
$x'=(x-\mu)/\sigma$ con statistiche del solo train; lo scaler resta salvato
nell'artifact e \file{metadati.json} registra gli hash sha256 dei dati. Il
generatore (\file{training/synthetic\_generator.py}) campiona profili clinici
da gaussiane troncate con vincolo sistolica $>$ diastolica ($\ge$ 10 mmHg) e
scarta i campioni troppo vicini alle soglie (margini 2.0/2.0/1.0/0.2/1.0/2.0;
oltre 52\,000 scartati per 40\,000 di train). Gli split sono congelati in
\file{data/processed/} (Tab.~\ref{tab:split}).

Il teacher (\codice{classifica\_rischio} in \file{safety\_rules.py}, soglie
critiche in Tab.~\ref{tab:soglie}) controlla validità, assegna gravità per parametro, cerca pattern
combinati (es.\ ipotensione + tachicardia) e decide: pattern/anomalia critica $\to$ \codice{alto}, altra anomalia $\to$
\codice{medio}, altrimenti \codice{basso}. Il teacher offline
(\codice{TeacherRules}) riusa la stessa funzione: label ed gate non possono
divergere (guard test anti-duplicazione) e l'accordo sulle proprie label è
1.0 su 20\,000, limite superiore dell'MLP. Resta rigido: soglie a gradino
manuali, punto di partenza non oracolo clinico.

% =====================================================================
\section{Modello MLP numpy e training}\label{sec:modello}

MLP 6~$\to$~32~$\to$~3 da 323 parametri ($(6{\times}32{+}32)+(32{\times}3{+}3)$,
inizializzazione di He): combinazioni lineari + ReLU nascosta e softmax in
uscita, allenato su cross-entropy con backpropagation; il
gradient check dà errore relativo $\sim$$10^{-10}$ (soglia test
$<2{\times}10^{-5}$). Early stopping a pazienza 5 su loss di validazione; la
griglia (hidden 16/32/64 $\times$ lr 0.01/0.05, batch 32/64, 25 epoche) sceglie
\codice{hidden=32, lr=0.05, batch=32} (loss val $\approx 0.051$) senza mai
guardare il test; seguono riaddestramento, artifact e report (in
\file{data/models/}: \file{mlp\_telemedicina.npz} e \file{report.json}).

Il legame teacher--studente è knowledge distillation (Hinton et al., 2015):
l'MLP approssima le regole con confini morbidi vedendo solo risposte del
teacher su sintetici, quindi ogni metrica è fedeltà alle regole; i 298 errori
di approssimazione sono contenuti dal gate.

% =====================================================================
\section{Sicurezza, incertezza e fallback}\label{sec:sicurezza}

Le prime righe di \codice{SupervisedAgent.predict} applicano le regole prima
di ogni inferenza: se dicono \codice{alto}, l'MLP non viene consultato e il
recall di sistema sui critici vale 1.0 per costruzione; la regola
\codice{max(rule-based, MLP)} permette al modello solo di alzare, mai di
abbassare (qualche falso positivo in più, accettato). Sui non critici la
confidenza $\max(p)$ sotto \codice{SOGLIA\_INCERTEZZA = 0.6} rimanda alle
regole con motivo registrato: la softmax non è calibrata, ma dove il modello
è debole il sistema diventa conservativo. Ogni analisi è tracciabile in
SQLite (\codice{modello\_usato}, \codice{classe\_mlp},
\codice{classe\_regola}, \codice{confidenza}, \codice{fallback},
\codice{motivo\_fallback}, \codice{override\_sicurezza}); se il DB non è
scrivibile \codice{analizza()} solleva \codice{ValueError}. Critici e invalidi
producono doppia notifica (riga \codice{notifiche} + stderr), emessa dal
service.

% =====================================================================
\section{Uso e conclusioni}\label{sec:uso}

Uso da root con bash + python3 tramite \file{progetto.sh}:

\begin{lstlisting}[caption={Comandi principali di \file{progetto.sh}.}]
./progetto.sh test              # suite (139 test) + regressione legacy (14)
./progetto.sh train             # rigenera artifact MLP + report.json
./progetto.sh analisi '{...}'   # analizza un paziente, persiste su SQLite
./progetto.sh dashboard         # dashboard statica in data/dashboard.html
./progetto.sh serve [porta]     # dashboard + API su 127.0.0.1:porta
\end{lstlisting}
Il training esplicito sovrascrive artifact e report in qualche minuto; se
l'artifact manca, il runtime usa il rule-based e lo dichiara. La dashboard
(unico punto con matplotlib, dipendenza di sviluppo) mostra in sola lettura
analisi, flusso, confronto teacher-vs-MLP, training e incertezza. Un caso come
\codice{(185,110,80,36.8,98,95)} (crisi ipertensiva) è intercettato dal gate,
che avvisa il medico.

Il sistema copre il task end-to-end (parametri $\to$ classe $\to$ messaggio
$\to$ alert $\to$ DB): l'MLP riproduce il teacher nel 98.5\,\% dei casi, il
gate porta il recall sui critici a 1.0, la soglia 0.6 dichiara l'incertezza
invece di nasconderla. Limiti: dataset sintetico, test ottimistico (0.9851
vs 0.9324), soglie legacy non validate da clinici, softmax non calibrata,
notifica persa se il suo insert fallisce dopo quello dell'analisi,
thread-safety non garantita. Sviluppi: dati reali anonimizzati, calibrazione,
SHAP/LIME, drift monitoring, più classi o indice continuo.

% =====================================================================
\appendix
\section{Tabelle, esempio e PEAS}\label{sec:appendice}

\begin{table}[htbp]
\centering
\caption{Messaggi al paziente e allerta medico per classe (da \file{services/analysis\_service.py}).}
\label{tab:messaggi}
\small
\begin{tabular}{lp{8cm}c}
\toprule
\textbf{Classe} & \textbf{Messaggio al paziente} & \textbf{Allerta} \\
\midrule
\codice{basso} & ``I tuoi parametri vitali risultano nella norma. Continua il monitoraggio regolare.'' & no \\
\codice{medio} & ``Sono state rilevate alcune anomalie. Ti consigliamo di contattare il tuo medico e ripetere la misurazione.'' & no \\
\codice{alto} & ``Sono stati rilevati parametri critici. Ti consigliamo di contattare subito un medico o recarti al pronto soccorso.'' & sì \\
\codice{errore} & messaggio di errore con la lista dei problemi rilevati & sì \\
\bottomrule
\end{tabular}
\end{table}

\begin{table}[htbp]
\centering
\caption{Le sei feature: unità e range normale (da \codice{RANGE\_NORMALI}).}
\label{tab:feature}
\small
\begin{tabular}{lll}
\toprule
\textbf{Feature} & \textbf{Unità} & \textbf{Range normale} \\
\midrule
\codice{pressione\_sistolica} & mmHg & 90--140 \\
\codice{pressione\_diastolica} & mmHg & 60--90 \\
\codice{frequenza\_cardiaca} & bpm & 60--100 \\
\codice{temperatura} & \textdegree{}C & 36.0--37.5 \\
\codice{saturazione\_ossigeno} & \% & 95--100 \\
\codice{glicemia} & mg/dL & 70--140 \\
\bottomrule
\end{tabular}
\end{table}

\begin{table}[htbp]
\centering
\caption{Soglie critiche per parametro (da \file{safety\_rules.py}).}
\label{tab:soglie}
\small
\begin{tabular}{ll}
\toprule
\textbf{Parametro} & \textbf{Soglia critica} \\
\midrule
\codice{pressione\_sistolica} & $\geq$ 180 mmHg \\
\codice{pressione\_diastolica} & $\geq$ 110 mmHg \\
\codice{frequenza\_cardiaca} & $\geq$ 120 oppure $\leq$ 50 bpm \\
\codice{temperatura} & $\geq$ 38.5 oppure $\leq$ 35.0 \textdegree{}C \\
\codice{saturazione\_ossigeno} & $\leq$ 90 \% \\
\codice{glicemia} & $\geq$ 200 oppure $\leq$ 60 mg/dL \\
\bottomrule
\end{tabular}
\end{table}

\begin{table}[htbp]
\centering
\caption{Split congelato (seed base 41).}
\label{tab:split}
\small
\begin{tabular}{lrr}
\toprule
\textbf{Blocco} & \textbf{Seed} & \textbf{Campioni} \\
\midrule
train & 41 & 40\,000 \\
validation & 42 & 10\,000 \\
test & 43 & 20\,000 \\
\bottomrule
\end{tabular}
\end{table}

\begin{lstlisting}[caption={Esempio di analisi di un paziente stabile.}]
./progetto.sh analisi '{"pressione_sistolica":120,"pressione_diastolica":80,\
"frequenza_cardiaca":75,"temperatura":36.8,\
"saturazione_ossigeno":98,"glicemia":95}'
\end{lstlisting}
L'output riporta classe (\codice{basso}), messaggio al paziente e metadati
dell'esito (chi ha deciso, confidenza, eventuale fallback).

In termini PEAS: performance = accuratezza, correttezza delle notifiche e
tempi di risposta; ambiente = paziente, wearable, medico e storico
(parzialmente osservabile, multi-agente, non deterministico, sequenziale,
dinamico); attuatori = messaggio, notifica e scrittura DB; sensori = i sei
vitali --- per questo il progetto ibrida MLP e regole con precedenza delle
regole.

\begin{thebibliography}{9}
\bibitem{task} Task del progetto, \file{projectDescription/task.txt} (HW\_persia\_privato).
\bibitem{contract} Contratto di dominio: feature, classi, safety (fonte unica: \file{safety\_rules.py}).
\bibitem{impl} Mappa concetti del corso $\to$ progetto, \file{knowledge/09\_implicazioni\_progetto.md} (PEAS, supervised vs RL, neuro-simbolico, AI Act).
\bibitem{doc} Documentazione tecnica, \file{supervised\_telemedicina/documentazione.md} (pipeline, DB, design decisions, report seed 41).
\bibitem{goodfellow} I.~Goodfellow, Y.~Bengio, A.~Courville, \textit{Deep Learning}, MIT Press, 2016.
\bibitem{hinton} G.~Hinton, O.~Vinyals, J.~Dean, Distilling the Knowledge in a Neural Network, arXiv:1503.02531, 2015.
\bibitem{bishop} C.~M.~Bishop, \textit{Pattern Recognition and Machine Learning}, Springer, 2006.
\end{thebibliography}

\end{document}
